\chapter{Nilpotent Groups, Frattini Subgroups, and Fitting Subgroups}

\providecommand{\Syl}{\operatorname{Syl}}
\providecommand{\Max}{\operatorname{Max}}
\providecommand{\Fit}{\operatorname{Fit}}
\providecommand{\radn}{\operatorname{rad}}
\providecommand{\Core}{\operatorname{Core}}
\providecommand{\Soc}{\operatorname{Soc}}
\providecommand{\normal}{\trianglelefteq}
\providecommand{\nnormal}{\triangleleft}

\section{The Frattini Argument}

\begin{theorem}[Frattini Argument]
Let \(G\) be a finite group, and let \(N\normal G\). If
\[
    P\in \Syl_p(N),
\]
then
\[
    G=N_G(P)N.
\]
\end{theorem}

\begin{proof}
Since \(N\normal G\), conjugation by every element of \(G\) sends Sylow
\(p\)-subgroups of \(N\) to Sylow \(p\)-subgroups of \(N\). Hence \(G\) acts by
conjugation on \(\Syl_p(N)\).

By the Sylow theorem, \(N\) acts transitively on \(\Syl_p(N)\). Let \(g\in G\).
Then \(gPg^{-1}\in \Syl_p(N)\), so there exists \(n\in N\) such that
\[
    gPg^{-1}=nPn^{-1}.
\]
Thus
\[
    n^{-1}gP g^{-1}n=P,
\]
so \(n^{-1}g\in N_G(P)\). Hence \(g\in nN_G(P)\subseteq NN_G(P)\). Therefore
\[
    G=NN_G(P)=N_G(P)N,
\]
because \(N\normal G\).
\end{proof}

%----------------------------------------------------------------------------------------
\section{Nilpotent Groups}

\begin{definition}[Central Series and Nilpotent Groups]
A group \(G\) is called \textbf{nilpotent} if there exists a chain of subgroups
\[
    G=K_1>K_2>\cdots>K_s>K_{s+1}=1
\]
such that
\[
    K_i\normal G
\]
and
\[
    K_i/K_{i+1}\leq Z(G/K_{i+1})
\]
for every \(i\). Such a series is called a \textbf{central series}.
\end{definition}

Equivalently, the condition
\[
    K_i/K_{i+1}\leq Z(G/K_{i+1})
\]
means
\[
    [G,K_i]\leq K_{i+1}.
\]

\begin{proposition}
Every finite \(p\)-group is nilpotent.
\end{proposition}

\begin{proof}
Let \(G\) be a nontrivial finite \(p\)-group. Then
\[
    Z(G)\neq 1.
\]
Define the upper central series
\[
    1=Z_0(G)\leq Z_1(G)\leq Z_2(G)\leq \cdots
\]
by
\[
    Z_{i+1}(G)/Z_i(G)=Z(G/Z_i(G)).
\]
If \(Z_i(G)<G\), then \(G/Z_i(G)\) is a nontrivial \(p\)-group, so it has
nontrivial center. Hence \(Z_{i+1}(G)>Z_i(G)\). Since \(G\) is finite, this
process must eventually reach \(G\). Therefore \(G\) has a central series and
is nilpotent.
\end{proof}

\begin{proposition}
Subgroups and quotient groups of nilpotent groups are nilpotent.
\end{proposition}

\begin{proof}
Let
\[
    G=K_1>K_2>\cdots>K_s>K_{s+1}=1
\]
be a central series for \(G\).

If \(H\leq G\), then
\[
    H=H\cap K_1\geq H\cap K_2\geq \cdots \geq H\cap K_{s+1}=1
\]
is a central series for \(H\), after deleting repetitions. Indeed,
\[
    [H,H\cap K_i]\leq [G,K_i]\leq K_{i+1},
\]
so
\[
    [H,H\cap K_i]\leq H\cap K_{i+1}.
\]

If \(N\normal G\), then
\[
    G/N=K_1N/N\geq K_2N/N\geq \cdots \geq K_{s+1}N/N=N/N
\]
is a central series for \(G/N\), again after deleting repetitions.
\end{proof}

\begin{proposition}
A finite direct product of nilpotent groups is nilpotent.
\end{proposition}

\begin{proof}
It is enough to consider \(G=A\times B\). If
\[
    A=A_1>\cdots>A_{r+1}=1
\]
and
\[
    B=B_1>\cdots>B_{t+1}=1
\]
are central series, then one obtains a central series for \(A\times B\) by
first descending through the \(A_i\times B\), and then through the
\(1\times B_j\). The centrality condition follows componentwise.
\end{proof}

\begin{example}
Let
\[
    G=P\times Q,
\]
where \(P\) is a \(p\)-group, \(Q\) is a \(q\)-group, and \(p\neq q\). Since
finite \(p\)-groups and finite \(q\)-groups are nilpotent, \(G\) is nilpotent.
Moreover
\[
    Z(G)=Z(P)\times Z(Q),
\]
and
\[
    G/Z(G)\cong P/Z(P)\times Q/Z(Q).
\]
Repeating this construction gives a central series for \(G\).
\end{example}

\begin{example}[Non-nilpotent groups]
The following groups are not nilpotent.

\begin{enumerate}[label=(\arabic*)]
    \item \(A_4\cong C_2^2\rtimes C_3\). Its only nontrivial proper normal
    subgroup is
    \[
        V_4\cong C_2^2.
    \]
    If \(A_4\) were nilpotent, the last nontrivial term of a central series
    would be contained in \(Z(A_4)\). But
    \[
        Z(A_4)=1.
    \]
    Hence \(A_4\) is not nilpotent.

    \item \(S_4\). Its proper nontrivial normal subgroups are
    \[
        V_4\quad\text{and}\quad A_4.
    \]
    Since
    \[
        Z(S_4)=1,
    \]
    \(S_4\) cannot have a central series. Therefore \(S_4\) is not nilpotent.

    \item Let \(p>3\) be prime, and let
    \[
        D_{2p}=\langle a,b\mid a^p=b^2=1,\; bab^{-1}=a^{-1}\rangle.
    \]
    The only nontrivial proper normal subgroup is \(\langle a\rangle\), but
    \[
        \langle a\rangle\nleq Z(D_{2p}).
    \]
    Hence \(D_{2p}\) is not nilpotent.
\end{enumerate}
\end{example}

\begin{theorem}[Equivalent Conditions for Nilpotency]
For a finite group \(G\), the following are equivalent.

\begin{enumerate}[label=(\arabic*)]
    \item \(G\) is nilpotent.

    \item For every proper subgroup \(H<G\),
    \[
        H<N_G(H).
    \]

    \item Every maximal subgroup \(M\) of \(G\) is normal and has prime index.
    Equivalently,
    \[
        G/M\cong C_p
    \]
    for some prime \(p\).

    \item \(G\) is the direct product of its Sylow subgroups. Equivalently,
    every Sylow subgroup of \(G\) is normal.
\end{enumerate}
\end{theorem}

\begin{proof}
\((1)\Rightarrow (2)\).
Let
\[
    G=K_1>K_2>\cdots>K_s>K_{s+1}=1
\]
be a central series for \(G\), and let \(H<G\). Choose \(i\) such that
\[
    K_{i+1}\leq H
    \qquad\text{but}\qquad
    K_i\nleq H.
\]
Since
\[
    K_i/K_{i+1}\leq Z(G/K_{i+1}),
\]
we have
\[
    [H,K_i]\leq K_{i+1}\leq H.
\]
Therefore \(K_i\leq N_G(H)\). Since \(K_i\nleq H\), it follows that
\[
    H<N_G(H).
\]

\((2)\Rightarrow (3)\).
Let \(M\) be a maximal subgroup of \(G\). By (2),
\[
    M<N_G(M)\leq G.
\]
By maximality of \(M\), this forces
\[
    N_G(M)=G.
\]
Thus \(M\normal G\).

Since \(M\) is maximal, \(G/M\) has no proper nontrivial subgroup. Hence
\(G/M\) is cyclic of prime order, say
\[
    G/M\cong C_p.
\]

\((3)\Rightarrow (4)\).
Let \(P\in \Syl_p(G)\). Then
\[
    P\leq N_G(P)\leq G.
\]
Suppose \(N_G(P)\neq G\). Choose a maximal subgroup \(M\) such that
\[
    N_G(P)\leq M<G.
\]
By (3), \(M\normal G\). Since \(P\leq M\) and \(P\) is a Sylow \(p\)-subgroup
of \(G\), \(P\) is also a Sylow \(p\)-subgroup of \(M\). By the Frattini
argument applied to \(M\normal G\),
\[
    G=M N_G(P).
\]
But \(N_G(P)\leq M\), so
\[
    G=M,
\]
a contradiction. Therefore
\[
    N_G(P)=G,
\]
so \(P\normal G\). Thus every Sylow subgroup of \(G\) is normal.

If the Sylow subgroups are
\[
    P_1,\dots,P_t,
\]
then they commute pairwise because
\[
    [P_i,P_j]\leq P_i\cap P_j=1
    \qquad (i\neq j).
\]
Hence
\[
    G=P_1\times \cdots \times P_t.
\]

\((4)\Rightarrow (1)\).
Each Sylow subgroup is a finite \(p\)-group, hence nilpotent. A finite direct
product of nilpotent groups is nilpotent. Therefore \(G\) is nilpotent.
\end{proof}

\begin{example}
Every finite abelian group is nilpotent.
\end{example}

\begin{example}
The group
\[
    A_4\cong C_2^2\rtimes C_3
\]
is not nilpotent. Indeed, a subgroup \(C_3\leq A_4\) is maximal, but it has
index \(4\), which is not prime, and it is not normal in \(A_4\).
\end{example}

%----------------------------------------------------------------------------------------
\section{The Frattini Subgroup}

\begin{definition}[Frattini Subgroup]
Let \(G\) be a finite group. Denote by
\[
    \Max(G)=\{M\mid M<G,\; M\text{ maximal}\}
\]
the set of maximal subgroups of \(G\). The \textbf{Frattini subgroup} of \(G\) is
\[
    \Phi(G)=\bigcap_{M\in \Max(G)}M.
\]
For the trivial group, we use the convention
\[
    \Phi(1)=1.
\]
\end{definition}

\begin{proposition}
The Frattini subgroup \(\Phi(G)\) is characteristic in \(G\).
\end{proposition}

\begin{proof}
Every automorphism of \(G\) permutes the maximal subgroups of \(G\). Hence it
fixes their intersection. Therefore
\[
    \Phi(G)  \operatorname{Char}(G).
\]
\end{proof}

\begin{corollary}
If \(N\normal G\), then
\[
    \Phi(N)\normal G.
\]
\end{corollary}

\begin{proof}
Since \(\Phi(N)  \operatorname{Char}(N)\) and \(N\normal G\), we have
\[
    \Phi(N)\normal G.
\]
\end{proof}

\begin{definition}[Nongenerator]
An element \(x\in G\) is called a \textbf{nongenerator} of \(G\) if, for every
subset \(S\subseteq G\),
\[
    \langle x,S\rangle=G
    \quad\Longrightarrow\quad
    \langle S\rangle=G.
\]
\end{definition}

\begin{lemma}[Characterization of the Frattini Subgroup]
The Frattini subgroup \(\Phi(G)\) consists exactly of the nongenerators of \(G\).
\end{lemma}

\begin{proof}
First suppose \(x\in \Phi(G)\). Assume
\[
    \langle x,S\rangle=G.
\]
If \(\langle S\rangle\neq G\), then \(\langle S\rangle\) is contained in some
maximal subgroup \(M<G\). Since \(x\in \Phi(G)\), we have \(x\in M\). Hence
\[
    \langle x,S\rangle\leq M<G,
\]
contradicting \(\langle x,S\rangle=G\). Therefore \(\langle S\rangle=G\), so
\(x\) is a nongenerator.

Conversely, suppose \(x\) is a nongenerator. If \(x\notin \Phi(G)\), then
there is a maximal subgroup \(M<G\) such that \(x\notin M\). By maximality,
\[
    \langle x,M\rangle=G.
\]
Since \(x\) is a nongenerator, this implies
\[
    \langle M\rangle=G,
\]
which is impossible because \(M<G\). Thus \(x\in M\) for every maximal
subgroup \(M\), and hence
\[
    x\in \Phi(G).
\]
\end{proof}

\begin{corollary}
If \(H\leq G\) and
\[
    H\Phi(G)=G,
\]
then
\[
    H=G.
\]
\end{corollary}

\begin{proof}
If \(H<G\), then \(H\) is contained in some maximal subgroup \(M<G\). Since
\[
    \Phi(G)\leq M,
\]
we get
\[
    H\Phi(G)\leq M<G,
\]
contradicting \(H\Phi(G)=G\). Therefore \(H=G\).
\end{proof}

\begin{example}
Let
\[
    G=\langle a\rangle\cong C_{p^2}.
\]
Then
\[
    \Phi(G)=\langle a^p\rangle.
\]
Thus every element of \(\langle a^p\rangle\) is a nongenerator.
\end{example}

\begin{example}
For the quaternion group \(Q_8\),
\[
    \Phi(Q_8)=Z(Q_8)=\{\pm 1\}.
\]
Thus the nongenerators of \(Q_8\) have order \(1\) or \(2\).
\end{example}

\begin{example}
Let
\[
    D_{2n}=\langle r,s\mid r^n=s^2=1,\; srs=r^{-1}\rangle
\]
be the dihedral group of order \(2n\). Then
\[
    \Phi(D_{2n})=\langle r^{\radn(n)}\rangle,
\]
where \(\radn(n)\) is the product of the distinct prime divisors of \(n\).
\end{example}

\begin{example}
If \(G\) is a finite nonabelian simple group, then
\[
    \Phi(G)=1.
\]
Indeed, \(\Phi(G)\normal G\), and \(\Phi(G)\neq G\) because \(G\) has maximal
subgroups. Thus \(\Phi(G)=1\). Hence no nonidentity element of \(G\) is a
nongenerator.

A stronger theorem says that for every nonidentity \(a\in G\), there exists
\(b\in G\) such that
\[
    \langle a,b\rangle=G.
\]
\end{example}

%----------------------------------------------------------------------------------------
\section{Frattini Subgroups and Nilpotency}

\begin{theorem}[Gasch\"utz]
Let \(N\normal G\) and suppose
\[
    N\leq \Phi(G).
\]
Then
\[
    G\text{ is nilpotent}
    \quad\Longleftrightarrow\quad
    G/N\text{ is nilpotent}.
\]
\end{theorem}

\begin{proof}
If \(G\) is nilpotent, then every quotient of \(G\) is nilpotent, so \(G/N\) is
nilpotent.

Conversely, suppose \(G/N\) is nilpotent. Let \(P\in \Syl_p(G)\). We show that
\(P\normal G\).

Since \(N\normal G\), the subgroup
\[
    PN/N
\]
is a Sylow \(p\)-subgroup of \(G/N\). Because \(G/N\) is nilpotent, all of its
Sylow subgroups are normal. Hence
\[
    PN/N\normal G/N,
\]
so
\[
    PN\normal G.
\]
Now \(P\in \Syl_p(PN)\). Applying the Frattini argument to \(PN\normal G\), we
obtain
\[
    G=N_G(P)PN.
\]
Since \(P\leq N_G(P)\), this becomes
\[
    G=N_G(P)N.
\]
But \(N\leq \Phi(G)\), so
\[
    N_G(P)\Phi(G)=G.
\]
By the previous corollary,
\[
    N_G(P)=G.
\]
Thus \(P\normal G\). Since this holds for every prime \(p\), all Sylow
subgroups of \(G\) are normal. Therefore \(G\) is nilpotent.
\end{proof}

\begin{corollary}
For a finite group \(G\), the following statements hold.

\begin{enumerate}[label=(\arabic*)]
    \item \(\Phi(G)\) is nilpotent.

    \item \(G\) is nilpotent if and only if \(G/\Phi(G)\) is nilpotent.

    \item \textup{(Wielandt)} \(G\) is nilpotent if and only if
    \[
        G'\leq \Phi(G).
    \]

    \item If \(G\) is a \(p\)-group, then
    \[
        G/\Phi(G)
    \]
    is elementary abelian.
\end{enumerate}
\end{corollary}

\begin{proof}
\((1)\).
Let
\[
    P\in \Syl_p(\Phi(G)).
\]
Since \(\Phi(G)\normal G\), the Frattini argument gives
\[
    G=\Phi(G)N_G(P).
\]
Because \(\Phi(G)\leq \Phi(G)\), the corollary above gives
\[
    N_G(P)=G.
\]
Hence \(P\normal G\), and in particular \(P\normal \Phi(G)\). Therefore every
Sylow subgroup of \(\Phi(G)\) is normal, so \(\Phi(G)\) is nilpotent.

\((2)\).
Apply Gasch\"utz's theorem with
\[
    N=\Phi(G).
\]

\((3)\).
Suppose \(G\) is nilpotent. Let \(M<G\) be maximal. Then \(M\normal G\), and
\[
    G/M\cong C_p
\]
for some prime \(p\). In particular, \(G/M\) is abelian, so
\[
    G'\leq M.
\]
Since this holds for every maximal subgroup \(M\), we get
\[
    G'\leq \Phi(G).
\]

Conversely, suppose
\[
    G'\leq \Phi(G).
\]
Then
\[
    (G/\Phi(G))'=1,
\]
so \(G/\Phi(G)\) is abelian, hence nilpotent. By (2), \(G\) is nilpotent.

\((4)\).
Let \(G\) be a \(p\)-group. By (3),
\[
    G'\leq \Phi(G),
\]
so \(G/\Phi(G)\) is abelian.

It remains to show that every element has order dividing \(p\). Let \(x\in G\).
For every maximal subgroup \(M<G\), we have
\[
    G/M\cong C_p.
\]
Hence
\[
    x^p\in M.
\]
Since this holds for every maximal subgroup \(M\),
\[
    x^p\in \Phi(G).
\]
Thus every element of \(G/\Phi(G)\) has order dividing \(p\). Therefore
\(G/\Phi(G)\) is elementary abelian.
\end{proof}

\begin{example}
Let
\[
    G=C_{p^2}\rtimes C_q=\langle a\rangle\rtimes \langle b\rangle,
\]
where \(p\neq q\) are primes, \(|a|=p^2\), and \(|b|=q\). Write
\[
    b^{-1}ab=a^u
\]
for some
\[
    u\in (\Z/p^2\Z)^\times
\]
satisfying
\[
    u^q\equiv 1\pmod{p^2}.
\]
Then
\[
    a^p b=ba^p
\]
if and only if
\[
    u\equiv 1\pmod p.
\]
But the kernel of
\[
    (\Z/p^2\Z)^\times\longrightarrow (\Z/p\Z)^\times
\]
has order \(p\). Since \(q\neq p\), the condition \(u^q\equiv 1\pmod{p^2}\) and
\(u\equiv 1\pmod p\) forces
\[
    u\equiv 1\pmod{p^2}.
\]
Therefore
\[
    a^p b=ba^p
    \quad\Longleftrightarrow\quad
    ab=ba.
\]
Equivalently,
\[
    G/\langle a^p\rangle\text{ is abelian}
    \quad\Longleftrightarrow\quad
    G\text{ is abelian}.
\]
\end{example}

%----------------------------------------------------------------------------------------
\section{The Fitting Subgroup}

\begin{definition}
Let \(p\) be a prime. The subgroup \(O_p(G)\) is defined to be the largest
normal \(p\)-subgroup of \(G\). Equivalently,
\[
    O_p(G)
    =
    \prod_{\substack{N\normal G\\ N\text{ a }p\text{-group}}}N
    =
    \bigcap_{P\in \Syl_p(G)}P.
\]
\end{definition}

\begin{proposition}
For every prime \(p\), the subgroup \(O_p(G)\) is characteristic in \(G\). Moreover,
\[
    O_p(G/O_p(G))=1.
\]
If \(p\neq q\), then
\[
    [O_p(G),O_q(G)]=1,
\]
and hence
\[
    O_p(G)O_q(G)=O_p(G)\times O_q(G).
\]
\end{proposition}

\begin{proof}
The subgroup \(O_p(G)\) is uniquely characterized as the largest normal
\(p\)-subgroup of \(G\), so it is characteristic.

Let \(\overline{G}=G/O_p(G)\). If \(O_p(\overline{G})\neq 1\), then its full
preimage in \(G\) is a normal \(p\)-subgroup strictly larger than \(O_p(G)\),
contradicting maximality. Hence
\[
    O_p(\overline{G})=1.
\]

Finally, if \(p\neq q\), then
\[
    [O_p(G),O_q(G)]\leq O_p(G)\cap O_q(G)=1.
\]
Therefore \(O_p(G)\) and \(O_q(G)\) centralize each other, and their product is
direct.
\end{proof}

\begin{lemma}[Fitting's Lemma]
If \(H\normal G\) and \(K\normal G\) are nilpotent normal subgroups of \(G\), then
\[
    HK
\]
is a nilpotent normal subgroup of \(G\).
\end{lemma}

\begin{proof}
Normality of \(HK\) is clear. It remains to prove nilpotency.

Replacing \(G\) by \(HK\), we may assume
\[
    G=HK.
\]
We argue by induction on the sum of the nilpotency classes of \(H\) and \(K\).

Since \(H\normal G\), its center \(Z(H)\) is characteristic in \(H\), hence normal
in \(G\). The quotient
\[
    G/Z(H)
\]
is generated by the nilpotent normal subgroups
\[
    H/Z(H)
    \qquad\text{and}\qquad
    KZ(H)/Z(H).
\]
The nilpotency class of \(H/Z(H)\) is smaller than that of \(H\), so by induction
\(G/Z(H)\) is nilpotent. Similarly, \(G/Z(K)\) is nilpotent.

Thus for sufficiently large \(m\),
\[
    \gamma_m(G)\leq Z(H)\cap Z(K),
\]
where \(\gamma_m(G)\) denotes the \(m\)-th term of the lower central series.
But
\[
    Z(H)\cap Z(K)\leq Z(G),
\]
because \(G=HK\). Hence
\[
    \gamma_{m+1}(G)=1.
\]
Therefore \(G\) is nilpotent.
\end{proof}

\begin{definition}[Fitting Subgroup]
The \textbf{Fitting subgroup} of \(G\), denoted
\[
    F(G),
\]
is the product of all normal nilpotent subgroups of \(G\). By Fitting's lemma,
\(F(G)\) is the unique largest nilpotent normal subgroup of \(G\).
\end{definition}

Since \(\Phi(G)\) is nilpotent and normal, we always have
\[
    \Phi(G)\leq F(G).
\]

\begin{theorem}
Let \(G\) be a finite solvable group, and let
\[
    |G|=p_1^{e_1}p_2^{e_2}\cdots p_t^{e_t}.
\]
Let
\[
    F=F(G).
\]
Then the following hold.

\begin{enumerate}[label=(\arabic*)]
    \item
    \[
        F=O_{p_1}(G)\times O_{p_2}(G)\times \cdots \times O_{p_t}(G).
    \]

    \item
    \[
        C_G(F)\leq F.
    \]
    In other words, \(F(G)\) is self-centralizing.

    \item Consequently,
    \[
        G/C_G(F)\hookrightarrow \Aut(F).
    \]
\end{enumerate}
\end{theorem}

\begin{proof}
\((1)\).
Let \(N\normal G\) be nilpotent. Then \(N\) is the direct product of its Sylow
subgroups:
\[
    N=Q_1\times Q_2\times \cdots \times Q_t,
\]
where \(Q_i\in \Syl_{p_i}(N)\). Since \(N\) is nilpotent, each \(Q_i\) is
characteristic in \(N\). Since \(N\normal G\), it follows that
\[
    Q_i\normal G.
\]
Thus
\[
    Q_i\leq O_{p_i}(G).
\]
Therefore every nilpotent normal subgroup \(N\) of \(G\) is contained in
\[
    O_{p_1}(G)\times \cdots \times O_{p_t}(G).
\]
Conversely, the product
\[
    O_{p_1}(G)\times \cdots \times O_{p_t}(G)
\]
is a nilpotent normal subgroup of \(G\). Hence it is contained in \(F(G)\).
This proves
\[
    F(G)=O_{p_1}(G)\times \cdots \times O_{p_t}(G).
\]

\((2)\).
Let
\[
    C=C_G(F).
\]
Since \(F \operatorname{Char}(G)\), the subgroup \(C\) is normal in \(G\). Also
\[
    C\cap F=Z(F).
\]
Set
\[
    W=C\cap F.
\]
Suppose, for contradiction, that
\[
    C>W.
\]
Then \(C/W\) is a nontrivial solvable normal subgroup of \(G/W\). Hence it has
a nontrivial abelian normal subgroup. In particular, there exists a subgroup
\(Y\) such that
\[
    W<Y\leq C,\qquad Y\normal G,
\]
and
\[
    Y/W
\]
is a nontrivial \(p\)-group for some prime \(p\).

Since \(Y\leq C=C_G(F)\), the subgroup \(Y\) centralizes \(F\), and hence
centralizes \(W\leq F\). Thus
\[
    W\leq Z(Y).
\]
Now \(Y/W\) is a \(p\)-group, hence nilpotent, and \(W\leq Z(Y)\) is nilpotent.
Therefore \(Y\) is nilpotent. Since \(Y\normal G\), maximality of \(F(G)\) gives
\[
    Y\leq F.
\]
But \(Y\leq C\), so
\[
    Y\leq C\cap F=W,
\]
contradicting \(W<Y\). Hence
\[
    C=W\leq F.
\]
Thus
\[
    C_G(F)\leq F.
\]

\((3)\).
The conjugation action of \(G\) on \(F\) gives a homomorphism
\[
    G\longrightarrow \Aut(F).
\]
Its kernel is precisely
\[
    C_G(F).
\]
Therefore
\[
    G/C_G(F)\hookrightarrow \Aut(F).
\]
\end{proof}

\begin{example}[Groups with Cyclic Sylow Subgroups]
Assume \(G\) is a finite group whose Sylow subgroups are all cyclic. A standard
theorem says that such a group is solvable.

Let
\[
    F=F(G).
\]
Since each \(O_p(G)\) is a subgroup of a cyclic Sylow \(p\)-subgroup, each
\(O_p(G)\) is cyclic. Therefore
\[
    F=\prod_{p\mid |G|}O_p(G)
\]
is cyclic.

By the preceding theorem,
\[
    C_G(F)\leq F.
\]
Since \(F\) is abelian, we also have
\[
    F\leq C_G(F).
\]
Thus
\[
    C_G(F)=F.
\]
Hence
\[
    G/F\hookrightarrow \Aut(F).
\]
Because \(F\) is cyclic, \(\Aut(F)\) is abelian. Therefore \(G/F\) is abelian.

Moreover, the Sylow subgroups of \(G/F\) are cyclic. A finite abelian group whose
Sylow subgroups are cyclic is cyclic. Hence \(G/F\) is cyclic.

Thus \(G\) is metacyclic: there exists a cyclic normal subgroup \(F\normal G\)
such that
\[
    G/F
\]
is cyclic.
\end{example}

